\corrigesChapitre{8}{Probabilités conditionnelles et Bayes}
\begin{corrige}{1}{Lire un conditionnement}
\begin{questions}
\item $0{,}4\times0{,}7=0{,}28$.
\item $1-0{,}7=0{,}3$.
\item Non : il manque la probabilité de $B$ dans le groupe $\bar A$, ou une information équivalente.
\end{questions}
\end{corrige}
\begin{corrige}{2}{Compléter un arbre}
\begin{questions}
\item $P(\bar A)=0{,}7$, $P_A(\bar B)=0{,}2$, $P_{\bar A}(\bar B)=0{,}8$.
\item $0{,}3\times0{,}8+0{,}7\times0{,}2=0{,}38$.
\item $0{,}7\times0{,}2=0{,}14$.
\end{questions}
\end{corrige}
\begin{corrige}{3}{Indépendance}
\begin{questions}
\item Oui : $0{,}2=0{,}5\times0{,}4$.
\item $0{,}2/0{,}5=0{,}4=P(B)$.
\item Non : leur intersection a une probabilité positive. Indépendance et incompatibilité sont deux notions différentes.
\end{questions}
\end{corrige}
\begin{corrige}{4}{Deux urnes}
\begin{questions}
\item $\frac13\frac34+\frac23\frac14=5/12$.
\item $P_R(U)=(1/4)/(5/12)=3/5$.
\item $3/5>1/3$ : observer une boule rouge rend l’urne $U$ plus probable.
\end{questions}
\end{corrige}
\begin{corrige}{5}{Production et défauts}
\begin{questions}
\item $0{,}6\times0{,}02+0{,}4\times0{,}05=0{,}032$, soit $3{,}2\,\%$.
\item $0{,}4\times0{,}05/0{,}032=5/8=62{,}5\,\%$.
\item Les deux machines n’ont pas le même poids dans la production ; il faut une moyenne pondérée.
\end{questions}
\end{corrige}
\begin{corrige}{6}{Un tableau de test}
\begin{questions}
\item $s=45/50=0{,}9$ ; $c=931/950=0{,}98$.
\item $45/(45+19)=45/64\simeq70{,}31\,\%$.
\item $931/(931+5)=931/936\simeq99{,}47\,\%$.
\end{questions}
\end{corrige}
\begin{corrige}{7}{Trois origines}
\begin{questions}
\item $0{,}002+0{,}006+0{,}020=0{,}028$.
\item $0{,}020/0{,}028=5/7$.
\item Elles valent $1/14$, $3/14$ et $10/14$ ; leur somme vaut $1$.
\end{questions}
\end{corrige}
\begin{corrige}{8}{Retrouver une proportion}
\begin{questions}
\item $0{,}35=0{,}8p+0{,}2(1-p)$.
\item $0{,}15=0{,}6p$, donc $p=1/4$.
\item $0{,}8\times0{,}25/0{,}35=4/7\simeq57{,}14\,\%$.
\end{questions}
\end{corrige}
\begin{corrige}{9}{Prévalence et valeur prédictive}
\begin{questions}
\item $V(p)=0{,}9p/(0{,}85p+0{,}05)$.
\item Le dénominateur est positif : $1{,}8p>0{,}85p+0{,}05$, donc $p>1/19\simeq5{,}26\,\%$.
\item La fiabilité d’un résultat positif dépend de la proportion initiale de $M$, même si la sensibilité du test reste $90\,\%$.
\end{questions}
\end{corrige}
\begin{corrige}{10}{Deux résultats positifs}
\begin{questions}
\item $0{,}9^2=0{,}81$ sachant $M$ ; $0{,}1^2=0{,}01$ sachant $\bar M$.
\item $0{,}1\times0{,}81/(0{,}081+0{,}009)=0{,}9$.
\item Les produits conditionnels $0{,}9^2$ et $0{,}1^2$ ne seraient plus justifiés ; il faudrait les probabilités conjointes des deux résultats.
\end{questions}
\end{corrige}
\begin{corrige}{11}{L’erreur de l’inversion}
\begin{questions}
\item On peut construire un arbre avec ces trois poids ; tous sont entre $0$ et $1$.
\item $P(B)=0{,}08+0{,}18=0{,}26$ ; $P_B(A)=0{,}08/0{,}26=4/13\simeq0{,}3077$, tandis que $P_A(B)=0{,}8$.
\item Il faut utiliser Bayes et le taux a priori. Inverser les lettres sans recalcul n’est pas une règle.
\end{questions}
\end{corrige}
\begin{corrige}{12}{Une preuve de croissance}
\begin{questions}
\item Le dénominateur est $f+(s-f)p$ ; la règle du quotient donne $V\,\!^{\prime}(p)=sf/[f+(s-f)p]^2$.
\item Elle est strictement positive, donc $V$ est strictement croissante.
\item $V(0)=0$ et $V(1)=1$ : à performances conditionnelles fixées, la valeur prédictive dépend fortement de la proportion initiale.
\end{questions}
\end{corrige}
